\documentclass[11pt]{article}
\usepackage[margin=1in]{geometry}
\usepackage{amsmath,amssymb,amsthm}
\usepackage{graphicx}
\usepackage[hidelinks]{hyperref}

\newtheorem{theorem}{Theorem}
\newtheorem{lemma}{Lemma}
\newtheorem{corollary}{Corollary}

\newcommand{\N}{\mathbb N}
\newcommand{\F}{\mathcal F}
\newcommand{\Fr}{\mathcal F_{\rm Fr}}
\newcommand{\SB}{\mathsf{SB}}

\title{A Counterexample to the Filter-\(\pi\) FDD Question}
\author{fa\_banach\_001, agent\_lane\_16}
\date{June 28, 2026}

\begin{document}
\maketitle

\section*{Source Question}

The source paper is Tomasz Kania and Jaroslaw Swaczyna,
\emph{The \(\pi\)-property of a Banach space along a filter},
arXiv:2508.21242.  Question 1.2 asks:
\begin{quote}
Suppose that a separable Banach space has the \(\mathcal F\)-\(\pi\)-property
with respect to a certain filter \(\mathcal F\). Does \(X\) have a
finite-dimensional decomposition (FDD)?
\end{quote}

\begin{center}
\includegraphics[width=.95\linewidth]{figures/open_problem_crop.png}
\end{center}

\section*{Result}

\begin{theorem}
Every separable Banach space \(X\) has the \(\mathcal F\)-\(\pi\)-property
for some filter \(\mathcal F\) on \(\N\) containing the Frechet filter.
\end{theorem}

\begin{corollary}
Question 1.2 of arXiv:2508.21242 has a negative answer.  Indeed, Enflo's
separable Banach space without the approximation property has the
\(\mathcal F\)-\(\pi\)-property for some such filter \(\mathcal F\), but it
does not have an FDD.
\end{corollary}

\section*{Proof Intuition}

The filter is allowed to depend on the space and on the chosen sequence of
finite-rank projections.  A separable Banach space has enough finite-rank
projections to approximate the identity on any fixed finite set of vectors:
project onto the finite-dimensional span of those vectors.  Separability lets
us choose a countable strong-operator dense list of norm-bounded finite-rank
projections, and we repeat every projection in that list infinitely many
times.  The desired filter is then generated by the sets of indices on which
this sequence approximates a given vector.  Finite intersections stay infinite
because any finite list of vector constraints is handled by one finite-rank
projection, which appears infinitely often.

\section*{Proof}

\begin{lemma}
Let \(X\) be separable.  For each \(m\in\N\), the set
\[
\mathcal P_m=\{P\in \mathcal L(X): P^2=P,\ \dim P[X]<\infty,\ \|P\|\le m\}
\]
has a countable dense subset for the strong operator topology restricted to
\(\mathcal P_m\).
\end{lemma}

\begin{proof}
Fix a dense sequence \((x_j)\) in \(X\).  On the norm ball
\(\{T\in\mathcal L(X):\|T\|\le m\}\), the metric
\[
d(S,T)=\sum_{j=1}^\infty 2^{-j}\min\{1,\|(S-T)x_j\|\}
\]
induces the strong operator topology.  This ball embeds into the countable
product \(\prod_j m\|x_j\|B_X\), which is separable because \(X\) is separable.
Hence the ball is separable in this topology, and so is its subspace
\(\mathcal P_m\).
\end{proof}

For each \(m\), choose a countable strong-operator dense set
\(\mathcal D_m\subseteq\mathcal P_m\).  Let \((Q_n)\) be a sequence that
enumerates \(\bigcup_m\mathcal D_m\) and repeats each of its elements
infinitely many times.

We first record the finite approximation property of this sequence.  Given
vectors \(x_1,\ldots,x_r\in X\) and \(\varepsilon>0\), let
\(E=\operatorname{span}\{x_1,\ldots,x_r\}\).  Since \(E\) is finite
dimensional, it is complemented in \(X\): there is a bounded projection
\(P:X\to E\).  Pick \(m>\|P\|\).  By strong-operator density of
\(\mathcal D_m\), there is \(Q\in\mathcal D_m\) such that
\[
\max_{1\le i\le r}\|Qx_i-Px_i\|<\varepsilon.
\]
Since \(Px_i=x_i\), this says \(\max_i\|Qx_i-x_i\|<\varepsilon\).  Because
the projection \(Q\) occurs infinitely many times in \((Q_n)\), the set of
indices \(n\) satisfying this finite list of inequalities is infinite.

For \(x\in X\) and \(k\in\N\), set
\[
A(x,k)=\{n\in\N:\|Q_nx-x\|<1/k\}.
\]
The preceding paragraph implies that every finite intersection of sets
\(A(x,k)\), even after intersecting with an arbitrary cofinite subset of
\(\N\), is nonempty; in fact it is infinite.  Therefore the family consisting
of all cofinite subsets of \(\N\) together with all \(A(x,k)\) generates a
proper filter \(\mathcal F\) containing the Frechet filter.

Now fix \(x\in X\) and \(\varepsilon>0\).  Choose \(k\) with \(1/k<\varepsilon\).
Since \(A(x,k)\in\mathcal F\), for all \(n\in A(x,k)\) we have
\(\|Q_nx-x\|<\varepsilon\).  Thus
\[
\lim_{n,\mathcal F}\|Q_nx-x\|=0.
\]
The sequence \((Q_n)\) consists of finite-rank projections, so it witnesses
the \(\mathcal F\)-\(\pi\)-property of \(X\).

For the corollary, let \(X\) be Enflo's separable Banach space without the
approximation property.  The theorem gives the \(\mathcal F\)-\(\pi\)-property
for some filter \(\mathcal F\) containing the Frechet filter.  If \(X\) had an
FDD \((F_n)\), the partial-sum projections \(S_N:X\to F_1+\cdots+F_N\) would
be finite-rank and would converge pointwise to the identity.  By the uniform
boundedness principle, \(\sup_N\|S_N\|<\infty\).  Hence \(X\) would have the
bounded approximation property, and therefore the approximation property,
contradicting Enflo's construction.  Thus this \(X\) is a counterexample to
Question 1.2.

Finally, if the intended reading of the question is restricted to analytic
filters, the same counterexample persists.  Theorem B of the source paper says
that whenever a sequence witnesses the \(\mathcal F\)-\(\pi\)-property for an
arbitrary filter \(\mathcal F\), there is an analytic filter
\(\mathcal G\subseteq\mathcal F\) for which the same sequence is still a
witness.  If needed to maintain the paper's convention that all filters contain
the Frechet filter, replace \(\mathcal G\) by the filter generated by
\(\mathcal G\cup\mathcal F_{\rm Fr}\).  This filter is still analytic, is still
contained in \(\mathcal F\), contains the Frechet filter, and still contains the
sets that witness convergence.

\section*{Verification Notes}

The filter contains the Frechet filter by construction.  The proof does not
use any approximation property of \(X\); it uses only separability and the
elementary fact that finite-dimensional subspaces of Banach spaces are
complemented.  The resulting sequence need not be uniformly bounded.  This is
the exact point where the construction escapes the source paper's observation
that a uniformly bounded \(\mathcal F\)-\(\pi\)-basis contains an ordinary
\(\pi\)-basis subsequence.

This result does not address the usual Frechet-filter \(\pi\)-property, nor
does it answer whether the classical \(\pi\)-property implies an FDD.

\section*{Novelty Check}

The local run indexes
\texttt{registry\_index.tsv}, \texttt{solutions/index.tsv},
\texttt{attempts/index.tsv}, and \texttt{proof\_gaps/index.tsv} were searched
for \texttt{2508.21242}, the source title, \texttt{FDD},
\texttt{finite-dimensional decomposition}, and close filter-\(\pi\) phrases;
no prior packet for this target was found.  A bounded external web search on
June 28, 2026 using the exact title and phrases such as
\texttt{"every separable Banach space" "filter" "pi-property"
"finite-rank projections"} found the source arXiv page but no existing
statement of this universal counterexample route.  This is not a comprehensive
MathSciNet/Zentralblatt search.

\begin{thebibliography}{9}
\bibitem{KS}
T. Kania and J. Swaczyna,
\emph{The \(\pi\)-property of a Banach space along a filter},
arXiv:2508.21242, 2025.

\bibitem{Enflo}
P. Enflo,
\emph{A counterexample to the approximation problem in Banach spaces},
Acta Mathematica 130 (1973), 309--317.
\end{thebibliography}

\end{document}
